\documentclass[a4paper,11pt]{article}
\usepackage[utf8]{inputenc}
\usepackage[T1]{fontenc}
\usepackage[french]{babel}
\usepackage{graphicx}
\usepackage{amsmath}
\usepackage{amsthm}
\theoremstyle{definition}
\usepackage{geometry}
\usepackage{fancyhdr}
\usepackage{hyperref}
\usepackage[usenames,dvipsnames,svgnames,table]{xcolor}
\usepackage{enumerate}
\usepackage{tcolorbox}
\tcbuselibrary{skins,breakable}
\usepackage{tikz}
\usepackage{pgfplots}
\pgfplotsset{compat=1.18}
\usetikzlibrary{arrows.meta}

\geometry{a4paper, scale=0.8}

%% ─── Header / footer ─────────────────────────────────────────────
\pagestyle{fancy}
\fancyhf{}
\fancyhead[L]{2025-2026}
\fancyhead[C]{\textbf{\courseName}}
\fancyhead[R]{\large GEII -- BUT2}
\fancyfoot[C]{\thepage}
\fancyfoot[L]{\includegraphics[height=3em]{../../common/logos_couleurs/IUTA.png}}
\fancyfoot[R]{\href{mailto:prenom.nom@univ-smb.fr}{prenom.nom@univ-smb.fr}}

%% ─── Course metadata ─────────────────────────────────────────────
\newcommand{\courseName}{Outils Mathématiques et Logiciel --- OML3}
\newcommand{\chapterTitle}{TD 1 -- Série de Fourier}

%% ─── Custom title ────────────────────────────────────────────────
\makeatletter
\renewcommand\maketitle{{\raggedright\begin{center}
    {\Huge\bfseries\@title}
\end{center}}}
\makeatother

%% ─── Solution toggle ─────────────────────────────────────────────
%% Compile with \showsolutionstrue for the teacher version.
\newif\ifshowsolutions
\showsolutionsfalse

\ifshowsolutions
    \newtcolorbox{solution}{
        colback=red!5!white, colframe=red!50!white,
        breakable, title=\textbf{Solution :}
    }
\else
    \NewEnviron{solution}{}{}
\fi

%% ─── Theorem / exercise environments ─────────────────────────────
\newtheorem{definition}{Définition}[section]
\newtheorem{proposition}{Propriété}[section]
\newtheorem{theorem}{Théorème}[section]
\newtheorem{exemple}{Exemple}[section]
\newtheorem{remark}{Remarque}[section]

%% Exercice : affiche "Exercice N (titre)." et referme aussitôt le
%% paragraphe avec \par. Volontairement PAS un \newtheorem : amsthm
%% reporte l'affichage du titre via \everypar, qui ne se déclenche
%% qu'au démarrage d'un paragraphe. Si le corps commence directement
%% par une liste (\begin{enumerate}), aucun paragraphe ne démarre
%% avant le premier \item, et le titre se retrouve collé au "1." de
%% la liste -- même avec une ligne vide dans la source. Cet
%% environnement affiche le titre immédiatement, donc ne dépend pas
%% de ce mécanisme et fonctionne quel que soit le contenu qui suit.
\newcounter{exercice}   % global counter — do not add [section]
\newenvironment{exercice}[1][]{%
  \refstepcounter{exercice}%
  \par\noindent\textbf{Exercice \theexercice}\if\relax\detokenize{#1}\relax\else\ (#1)\fi.\par\noindent\ignorespaces\vspace{1em}
}{\par\vspace{1em}}

%% ─── Section numbering ───────────────────────────────────────────
\renewcommand{\thesection}{{\arabic{section}}}

%% ─── Helpers ─────────────────────────────────────────────────────
\definecolor{niceblue}{RGB}{0,128,172}
\newcommand{\highlight}[1]{{\color{niceblue}\emph{#1}}}

\usepackage{amsmath,amsfonts,bm}

%% ─── Operators ───────────────────────────────────────────────────
\DeclareMathOperator*{\argmax}{arg\,max}
\DeclareMathOperator*{\argmin}{arg\,min}
\DeclareMathOperator{\MSE}{MSE}
\DeclareMathOperator{\sign}{sign}
\DeclareMathOperator{\sinc}{sinc}
\DeclareMathOperator{\tr}{tr}

%% ─── Number sets ─────────────────────────────────────────────────
\def\R{{\mathbb{R}}}
\def\C{{\mathbb{C}}}
\def\N{{\mathbb{N}}}
\def\Z{{\mathbb{Z}}}

%% ─── Infinitesimal and units ────────────────────────────────────
\def\d{{\rm d}}
\def\Hz{{\mbox{Hz}}}
\def\m{{\mbox{m}}}

%% ─── Complex / signal ───────────────────────────────────────────
\DeclareMathOperator{\complexj}{\mathsf{j}}   % imaginary unit

%% ─── Probability ─────────────────────────────────────────────────
\newcommand{\Esp}[2]{\mathbb{E}_{#1}\!\left[#2\right]}
\renewcommand{\Pr}{\mathbb{P}}

%% ─── Linear algebra ──────────────────────────────────────────────
\newcommand{\abs}[1]{\left|{#1}\right|}
\newcommand{\norm}[1]{\left\|#1\right\|}
\newcommand{\inner}[2]{\left\langle#1,\,#2\right\rangle}

%% ─── Logic ───────────────────────────────────────────────────────
\newcommand{\then}{\Rightarrow}
\newcommand{\iif}{\Leftrightarrow}

%% ─── Draft annotation ────────────────────────────────────────────
\newcommand{\red}[1]{{\color{red}#1}}

%% ─── Shared figure color (used by common/figures/*/*.tex bodies) ──
\definecolor{niceblue}{RGB}{0,128,172}

%% ─── Rappel header (styled subtitle for reference-card asides) ───
%% Requires \definecolor{niceblue}{...} to be defined by the including document.
\newcommand{\rappelhead}[1]{%
    \bigskip
    {\color{niceblue}\bfseries #1}\\[-0.4em]
    {\color{niceblue!35}\rule{\linewidth}{0.6pt}}
    \smallskip

}


\title{\chapterTitle}
\author{}
\date{}

%% ─────────────────────────────────────────────────────────────────
\begin{document}
\fontfamily{bch}\selectfont
\fontsize{12pt}{14pt}\selectfont

\maketitle
\thispagestyle{fancy}
\setcounter{exercice}{0}

\vspace{2em}
\textit{Tous les exercices ne seront pas abordés lors des séances de TD. Pour vos révisions il est recommandé de revoir précisément les 3 premiers exercices de chaque TD et de vous confronter au moins une fois à chaque exercice de la feuille de TD. Les corrections seront mises en ligne sur le Moodle du cours dans la semaine qui suit le TD.}
\vspace{2em}


%% ── Ex 1 : CLASSIFY / DEFINE ────────────────────────────────────
%% Purpose: activate vocabulary from the CM.
%% No computation harder than substituting into a definition.
%% At least one positive, one negative, one boundary case.
\begin{exercice}
    Pour les fonctions suivantes, indiquer si elles sont périodiques et si oui préciser leurs éventuelles périodes.
    \begin{enumerate}[\bfseries 1.]
        \item $f(t) = \cos(3t)$
        \begin{solution}
            $f$ est périodique : $f(t+T) = f(t) \iff \cos(3t+3T) = \cos(3t) \iff 3T = 2k\pi,\ k\in\Z$. La plus petite valeur positive est $T = \dfrac{2\pi}{3}$.
        \end{solution}

        \item $g(t) = t^2$
        \begin{solution}
            $g$ n'est pas périodique. Si elle l'était de période $T>0$, on aurait $\forall t\in\R,\ (t+T)^2 = t^2$, soit $2tT + T^2 = 0$ pour tout $t$ : impossible pour $T\neq 0$.
        \end{solution}

        \item $h(t) = \cos(2t) + \sin(3t)$
        \begin{solution}
            $\cos(2t)$ est périodique de période $\pi$, $\sin(3t)$ de période $\dfrac{2\pi}{3}$. La somme est périodique, de période le plus petit commun multiple des deux : $T = 2\pi$ (on vérifie $2\pi = 2\times\pi = 3\times\frac{2\pi}{3}$, et aucune valeur plus petite ne convient simultanément pour les deux termes).
        \end{solution}
    \end{enumerate}
\end{exercice}


%% ── Ex 2–3 : PROVE ───────────────────────────────────────────────
%% Purpose: make students re-derive a theorem stated in the CM.
%% Structure: easy direction → hard direction → verify on an example.


\begin{exercice}

    \begin{enumerate}[\bfseries 1.]
        \item
        Montrer que
        $$
        \cos^2(x) = \frac{\cos(2x) +1}{2}
        $$
        \begin{solution}
            D'après la formule d'addition, $\cos(2x) = \cos(x+x) = \cos^2(x) - \sin^2(x) = \cos^2(x) - (1-\cos^2(x)) = 2\cos^2(x) - 1$, d'où $\cos^2(x) = \dfrac{\cos(2x)+1}{2}$.
        \end{solution}
        \item Calculer
        $$
        \int_{0}^{2\pi} \cos^2(x) \d x
        $$
        \begin{solution}
            $$
                \int_0^{2\pi} \cos^2(x)\d x = \int_0^{2\pi} \frac{\cos(2x)+1}{2}\d x = \frac12\left[\frac{\sin(2x)}{2} + x\right]_0^{2\pi} = \frac12(0 + 2\pi - 0) = \pi.
            $$
        \end{solution}
    \end{enumerate}
\end{exercice}

\begin{exercice}[Calcul de valeurs moyennes]
    On rappelle que la valeur moyenne d'un signal $T$-périodique $f$ est définie par
    $$
        a_0 = \frac{1}{T}\int_0^T f(t)\,\d t.
    $$
   \begin{enumerate}[\bfseries 1.]
        \item Calculer la valeur moyenne de $f(t) = \cos(\omega t)$ ($\omega\neq 0$).
        \begin{solution}
            $$
                \bar{f} = \frac{1}{T}\int_0^T \cos(\omega t)\,\d t = \frac{1}{T}\left[\frac{\sin(\omega t)}{\omega}\right]_0^T = \frac{\sin(\omega T)}{\omega T} = 0
            $$
            car $\omega T = 2\pi$ (une période complète), donc $\sin(\omega T) = 0$. Plus généralement, la valeur moyenne d'un cosinus (ou d'un sinus) sur un nombre entier de périodes est toujours nulle.
        \end{solution}

        \item Calculer la valeur moyenne de $f_{\rm carr\acute{e}}$ (défini en cours, niveaux $-1$ et $1$). Généraliser au cas d'un signal carré prenant les valeurs $a$ et $b$ (chacune pendant une demi-période).
        \begin{solution}
            $$
                \bar{f}_{\rm carr\acute{e}} = \frac{1}{2\pi}\left(\int_0^\pi 1\,\d t + \int_\pi^{2\pi} (-1)\,\d t\right) = \frac{1}{2\pi}(\pi - \pi) = 0.
            $$
            Pour un signal prenant la valeur $a$ sur une demi-période et $b$ sur l'autre,
            $$
                \bar{f} = \frac{1}{T}\left(a\cdot\frac{T}{2} + b\cdot\frac{T}{2}\right) = \frac{a+b}{2}.
            $$
            On retrouve le cas précédent avec $a=1,b=-1$.
        \end{solution}

        \item Calculer la valeur moyenne de $f_{\rm scie}$ (défini en cours).
        \begin{solution}
            $$
                \bar{f}_{\rm scie} = \frac{1}{2\pi}\int_0^{2\pi} t\,\d t = \frac{1}{2\pi}\left[\frac{t^2}{2}\right]_0^{2\pi} = \frac{1}{2\pi}\cdot 2\pi^2 = \pi.
            $$
            C'est cohérent : une rampe linéaire de $0$ à $2\pi$ a pour moyenne le milieu, $\pi$.
        \end{solution}

        \item Un signal logique TTL vaut $0$~V ou $V_{cc}$, avec un rapport cyclique $\alpha\in[0,1]$ (fraction de la période passée à $V_{cc}$). Calculer sa valeur moyenne.
        \begin{solution}
            $$
                \bar{f} = \frac{1}{T}\left(\alpha T\cdot V_{cc} + (1-\alpha)T \cdot 0\right) = \alpha\, V_{cc}.
            $$
            \textbf{Conclusion :} c'est le principe de la commande par modulation de largeur d'impulsion (PWM/MLI) : faire varier $\alpha$ permet de régler la tension moyenne (donc la puissance transmise) sans dissiper d'énergie dans un composant résistif.
        \end{solution}
    \end{enumerate}
\end{exercice}


%% ── Ex 4–5 : APPLY ───────────────────────────────────────────────
%% Purpose: numerical problem with real or realistic GEII data.
%% Ends with an engineer's conclusion, not just a number.
\begin{exercice}[Décomposition en série de Fourier du signal carré]
    \begin{enumerate}[\bfseries 1.]
        \item En utilisant les formules de la décomposition en série de Fourier (forme réelle) vues en cours, calculer les coefficients $a_0, a_n, b_n$ jusqu'à l'ordre 5 de $f_{\rm carr\acute{e}}$.
        \begin{solution}
            Avec $T=2\pi$, $\omega_0 = 1$ :
            $$
                a_0 = \frac{1}{2\pi}\left(\int_0^\pi 1\,\d t - \int_\pi^{2\pi} 1\,\d t\right) = 0,
            $$
            ce qui est cohérent avec la valeur moyenne nulle calculée à l'exercice précédent.
            $$
                a_n = \frac{1}{\pi}\left(\int_0^\pi \cos(nt)\,\d t - \int_\pi^{2\pi}\cos(nt)\,\d t\right) = \frac{1}{\pi}\left(\frac{\sin(n\pi)}{n} - \frac{\sin(2n\pi)-\sin(n\pi)}{n}\right) = 0 \quad \forall n,
            $$
            puisque $\sin(n\pi) = \sin(2n\pi) = 0$ pour $n$ entier.
            $$
                b_n = \frac{1}{\pi}\left(\int_0^\pi \sin(nt)\,\d t - \int_\pi^{2\pi}\sin(nt)\,\d t\right) = \frac{2\left(1-(-1)^n\right)}{n\pi}.
            $$
            Donc $b_n = 0$ si $n$ est pair, et $b_n = \dfrac{4}{n\pi}$ si $n$ est impair. Jusqu'à l'ordre 5 :
            $$
                b_1 = \frac{4}{\pi},\quad b_2 = 0,\quad b_3 = \frac{4}{3\pi},\quad b_4 = 0,\quad b_5 = \frac{4}{5\pi}.
            $$
        \end{solution}

        \item Montrer, sans calcul d'intégrale, que $a_n = 0$ pour tout $n$.
        \begin{solution}
            $f_{\rm carr\acute{e}}$ est une fonction impaire (on vérifie $f_{\rm carr\acute{e}}(-t) = -f_{\rm carr\acute{e}}(t)$ sur une période). Le produit d'une fonction impaire par $\cos(nt)$ (paire) est impair. D'après l'Exercice~5 (invariance de l'intégrale par décalage d'une période), on peut calculer $a_n$ sur l'intervalle symétrique $[-\pi,\pi]$ plutôt que $[0,2\pi]$ ; l'intégrale d'une fonction impaire sur un intervalle symétrique étant nulle (Propriété~1.3 du cours), on obtient directement $a_n=0$.
        \end{solution}
    \end{enumerate}
\end{exercice}

\begin{exercice}
          Soit $f$ une fonction périodique et $T>0$ une période. 
            $$
            \text{Montrer que} \ \forall a\in \R \quad
            \int_a^{a+T} f(t) \d t = \int_{0}^{T} f(t) \d t.
            $$
        \begin{solution}
            Par la relation de Chasles,
            $$
                \int_a^{a+T} f(t)\d t = \int_a^{T} f(t) \d t + \int_T^{a+T} f(t) \d t.
            $$
            Dans la seconde intégrale, on pose $u = t - T$ (donc $\d u = \d t$) :
            $$
                \int_T^{a+T} f(t)\d t = \int_0^{a} f(u+T)\d u = \int_0^{a} f(u)\d u
            $$
            car $f$ est $T$-périodique. Ainsi,
            $$
                \int_a^{a+T} f(t)\d t = \int_a^{T} f(t)\d t + \int_0^{a} f(t)\d t = \int_0^{T} f(t)\d t
            $$
            en reconnaissant à nouveau la relation de Chasles.
        \end{solution}

\end{exercice}


%% (Preuve du corollaire de parité -- voir TD2, Exercice 3)

%% ── Ex 6 (optional) : EXTEND ─────────────────────────────────────
%% Purpose: harder result not fully covered in the CM (integration of
%% a Fourier series term by term) -- guided so students derive the
%% needed tool themselves (admis, question 2), then apply it.
%% Optional / approfondissement.
\begin{exercice}[Du signal carré au signal triangle --- approfondissement]
    On va retrouver la série de Fourier de $f_{\rm triangle}$ à partir de celle de $f_{\rm carr\acute{e}}$, déjà calculée à l'Exercice~4, plutôt qu'en recalculant les intégrales depuis la définition. Pour rappel,
    $$
    f_{\rm carr\acute{e}}(t) = \begin{cases} 1 & \text{pour} \ t \in [0,\frac{T}{2}[\\ -1 & \text{pour} \ t \in [\frac{T}{2},T[ \end{cases},
    $$
    $$
    f_{\rm triangle}(t) = \begin{cases} t & \text{pour} \ t \in [0,\frac{T}{2}]\\ T - t & \text{pour} \ t \in [\frac{T}{2},T] \end{cases}.
    $$
    On fixe $T=2\pi$ ($\omega_0=1$).

    \begin{enumerate}[\bfseries 1.]
        \item Calculer la dérivée de $f_{\rm triangle}$ sur $]0,\pi[$ puis sur $]\pi,2\pi[$. Que remarque-t-on ?
        \begin{solution}
            $f_{\rm triangle}(t)=t$ sur $[0,\pi]$ donc $f_{\rm triangle}'(t)=1$ sur $]0,\pi[$ ; $f_{\rm triangle}(t)=2\pi-t$ sur $[\pi,2\pi]$ donc $f_{\rm triangle}'(t)=-1$ sur $]\pi,2\pi[$ ($f_{\rm triangle}$ n'est pas dérivable aux points anguleux $t=0,\pi,2\pi$). On reconnaît exactement $f_{\rm carr\acute{e}}$ : $f_{\rm triangle}$ est une primitive de $f_{\rm carr\acute{e}}$.
        \end{solution}

        \item On admet que l'on peut intégrer terme à terme une série de Fourier : si $$g(t)=\sum_n^{+\infty} b_n\sin(nt),$$ une primitive de $g$ s'écrit
        $$
            G(t) = C - \sum_n^{+\infty} \frac{b_n}{n}\cos(nt), \qquad C\in\R.
        $$
        En utilisant les coefficients $b_n$ de $f_{\rm carr\acute{e}}$ (Exercice~4), en déduire qu'il existe $C\in\R$ tel que
        $$
            f_{\rm triangle}(t) = C - \sum_{n \text{ impair}} \frac{4}{n^2\pi}\cos(nt).
        $$
        \begin{solution}
            D'après l'Exercice~4, $f_{\rm carr\acute{e}}(t) = \displaystyle\sum_{n\text{ impair}} \frac{4}{n\pi}\sin(nt)$. En intégrant terme à terme (résultat admis) et d'après la question précédente,
            $$
                f_{\rm triangle}(t) = C - \sum_{n\text{ impair}} \frac{4}{n\pi\cdot n}\cos(nt) = C - \sum_{n\text{ impair}} \frac{4}{n^2\pi}\cos(nt).
            $$
        \end{solution}

        \item Déterminer $C$ en calculant directement la valeur moyenne de $f_{\rm triangle}$.
        \begin{solution}
            $$
                C = \bar{f}_{\rm triangle} = \frac{1}{2\pi}\left(\int_0^\pi t\,\d t + \int_\pi^{2\pi}(2\pi-t)\,\d t\right) = \frac{1}{2\pi}\left(\frac{\pi^2}{2}+\frac{\pi^2}{2}\right) = \frac{\pi}{2}.
            $$
        \end{solution}

        \item Conclure : écrire $f_{\rm triangle}(t)$ jusqu'à l'ordre 5, et vérifier que le résultat est cohérent avec le Corollaire~3.1 du cours.
        \begin{solution}
            $$
                f_{\rm triangle}(t) = \frac{\pi}{2} - \frac{4}{\pi}\cos(t) - \frac{4}{9\pi}\cos(3t) - \frac{4}{25\pi}\cos(5t) - \cdots
            $$
            La série ne contient que des cosinus, aucun terme en sinus : c'est cohérent avec le Corollaire~3.1, puisque $f_{\rm triangle}$ est une fonction paire.
        \end{solution}
    \end{enumerate}
\end{exercice}


%% (Impulsion périodique étroite -- voir TD2, Exercice 4)

\end{document}
